\begin{textsample}{POS Dim 1 – 2023\_c5 – Score 22.00 – 2023\_c5/enade\_2023\_eng\_comp\_e7383ab66f.txt}  \label{ex:f1_pos_004}
CAPÍTULO 4 DISTRIBUIÇÃO DE CONCEITOS A \textbf{finalidade} deste capítulo é apresentar um \textbf{panorama} nacional sobre a distribuição do Conceito Enade 2023.
Para traçar esse \textbf{panorama}, recorre -se a tabelas e gráficos, que são acompanhados de comentários e análises que contribuem para a contextualização dos dados neles expostos.
Para \textbf{visualizar} o \textbf{panorama} do Enade 2023, em nível nacional, na primeira seção deste capítulo, os dados sobre o número de cursos participantes são organizados por grande região.
Na \textbf{segunda}, o mapeamento do conceito do exame é feito por categoria administrativa das IES e modalidade de oferta dos cursos, de maneira \textbf{cruzada} às grandes regiões.
Na última seção, a lógica de apresentação dos dados é feita por organização acadêmica e, também, por grande região de funcionamento dos cursos, no país como um todo. 4.1 \textbf{PANORAMA} NACIONAL DA DISTRIBUIÇÃO DOS CONCEITOS A lógica adotada na apresentação dos dados, nesta seção e nas duas seguintes, pressupõe o registro de duas informações de natureza \textbf{reiterativa}.
O primeiro diz respeito aos valores possíveis do Enade, que variam em níveis de 1 a 5, sendo que, à medida que esses valores aumentam, o curso é mais bem avaliado, e ao fato de que aqueles cursos que tiveram somente um ou nenhum concluinte participante não obtêm o Conceito Enade, ficando sem conceito ( SC ), como ressaltado na apresentação deste relatório.
O segundo registro é que, caso haja diferença entre os cursos \textbf{tabulados} neste capítulo e no capítulo 1, ela corresponderá aos cursos que não foram avaliados, ou seja, àqueles, em princípio, sem estudantes concluintes que tivessem sido inscritos para a prova.
Na Tabela 4.1, são apresentadas a quantidade e a distribuição de cursos de Engenharia de Computação I participantes do Enade 2023, por faixa de conceito e grande região.
A partir dos dados da Tabela 4.1, nota -se que, dos 145 cursos participantes, 49 ( 33,8 \% ) classificaram -se com conceito 3.
Esse foi também o conceito modal nas seguintes regiões : Nordeste ( 44,0 \% ), Norte ( 45,5 \% ) e Sul ( 35,3 \% ).
O conceito 2 foi o segundo mais frequente em nível nacional ( 33,8 \%, correspondendo a 49 cursos ), tendo sido o segundo conceito mais frequente na maioria das regiões, exceto no Nordeste.
O conceito 4 foi o terceiro ( 16,6 \%, correspondendo a 24 cursos ).
Houve, ainda, 6 ( 4,1 \% ) cursos que receberam o conceito 5 e seis ( 4,1 \% ) cursos que receberam o conceito 1.
Dos 145 cursos de Engenharia de Computação I, 11 ( 7,6 \% ) ficaram sem conceito ( SC ).
A região Norte participou com 11 cursos, o que corresponde a 7,6 \% do total nacional.
Desses, cinco cursos receberam o conceito 3, o que \textbf{equivale} a 45,5 \% do total regional, o conceito modal na região, conforme mencionado.
A quatro cursos ( 36,4 \% ) atribuiu -se o conceito 2 ; a dois cursos ( 18,2 \% ), o conceito 4 ; e a nenhum curso ( 0,0 \% ) o conceito 5.
Ainda nessa região, nenhum curso ( 0,0 \% ) obteve conceito 1 e nenhum curso ( 0,0 \% ) ficou sem conceito ( SC ).
A região Nordeste participou com 25 cursos, correspondentes a 17,2 \% do total nacional.
Nessa região, 11 cursos ( 44,0 \% em termos regionais ) obtiveram conceito 3, o conceito modal para a região, como já comentado.
Os conceitos 4 e 2 foram atribuídos, respectivamente, a sete e a seis cursos ( 28,0 \% e 24,0 \%, respectivamente ).
Ainda nessa região, um curso ( 4,0 \% ) foi avaliado com o conceito 5 e nenhum curso ( 0,0 \% ) foi avaliado com o conceito 1.
Nenhum curso ( 0,0 \% ) ficou sem conceito ( SC ).
Para os 66 cursos participantes da região Sudeste ( 45,5 \% do total nacional ), o conceito 2 foi o conceito modal, como já anteriormente assinalado, tendo sido obtido por 28 cursos ( 42,4 \% ).
O conceito 3 foi atribuído a 19 cursos ( 28,8 \% ) e o conceito 4, a cinco cursos ( 7,6 \% ).
Receberam os conceitos 5 e 1, respectivamente, quatro cursos ( 6,1 \% ) e três cursos ( 4,5 \% ).
Nessa região, sete cursos ( 10,6 \% ) ficaram sem conceito ( SC ).
Os 34 cursos da região Sul corresponderam a 23,4 \% do total nacional.
Como já \textbf{indicado}, houve predominância do conceito 3 ( conceito modal ), atribuído a 35,3 \% dos cursos da região, ou seja, atribuído a 12 dos 34 cursos participantes na região Sul.
O conceito 4 foi atribuído a nove cursos ( 26,5 \% ) ; o conceito 2, seis cursos ( 17,6 \% ) ; e o conceito 1, a dois cursos ( 5,9 \% ).
Na região Sul, um curso ( 2,9 \% ) obteve conceito 5 e quatro cursos ( 11,8 \% ) ficaram sem conceito ( SC ).
Já Dos 9 cursos participantes na região Centro-Oeste ( 6,2 \% do total nacional ), 5 ( 55,6 \% em termos regionais ) receberam conceito 2, o conceito modal, como já destacado. dois cursos ( 22,2 \% ) obtiveram o conceito 3, um curso ( 11,1 \% ) obtiveram o conceito 4 e um curso ( 11,1 \% ), o conceito 1.
Nessa região, nenhum curso ( 0,0 \% ) obteve o conceito 5 e nenhum curso ( 0,0 \% ) ficou sem conceito ( SC ).
Em um gráfico de distribuição \textbf{cumulativa}, em cada valor de \textbf{abscissa} ( eixo x ), o valor apresentado no eixo das \textbf{ordenadas} ( eixo y ) é o valor acumulado da \textbf{variável} ( Conceito Enade ).
Em linhas gerais, em um gráfico desta mesma \textbf{tipologia}, \textbf{poligonais} mais à esquerda representam, no caso dos dados apresentados nesta seção, as grandes regiões com a pior distribuição, e \textbf{poligonais} mais à direita, as grandes regiões com os melhores conceitos.
A partir destas diretrizes, no Gráfico 4.1, é apresentada a distribuição \textbf{cumulativa} do Conceito Enade, segundo a grande região.
No Gráfico 4.1, observando a linha vertical \textbf{pontilhada}, que ajuda a \textbf{visualizar} o percentual de conceitos acumulados até o valor 3, ou seja, conceitos 1, 2 e 3, pode -se notar, por exemplo, que a região Centro-Oeste ( linha laranja ) apresenta 88,9 \% dos cursos com conceito 3 ou abaixo deste valor.
A região Sul ( linha \textbf{roxa} ), por outro lado, apresenta 66,7 \% dos cursos na mesma situação : com conceito 3 ou abaixo.
Destaca -se que, em algumas regiões, o conceito 5 não foi alcançado por nenhum curso, como ocorre nas regiões Norte ( linha azul ) e Centro-Oeste ( linha laranja ), nas quais todos os cursos têm conceito menor ou igual a 4, e as respectivas \textbf{poligonais} \textbf{cumulativas} alcançam 100 \% no conceito 4.
A região Centro-Oeste ( linha laranja ) apresenta os cursos com conceitos mais baixos, ao passo que a região Sul ( linha \textbf{roxa} ) apresenta a distribuição com valores maiores.
Os cursos em IES situadas na região Norte ( linha verde ) apresentam uma situação intermediária. 4.2 CONCEITOS POR CATEGORIA ADMINISTRATIVA, MODALIDADE DE OFERTA DOS CURSOS E GRANDE REGIÃO Na Tabela 4.2, é apresentada a distribuição dos cursos de Engenharia de Computação I participantes do Enade 2023, por categoria administrativa e modalidade de oferta, de acordo com os conceitos por eles alcançados, segundo as grandes regiões brasileiras.
De acordo com as informações da Tabela 4.2, dos 145 cursos participantes, 48 ( 33,1 \% ) eram \textbf{ministrados} em instituições públicas e 97 ( 66,9 \% ), em instituições privadas.
Quanto à modalidade de oferta, 126 ( 86,9 \% ) cursos eram presenciais e 19 ( 13,1 \% ), a distância.
De acordo com as informações da mesma tabela, em termos nacionais, entre instituições públicas e privadas, observa -se que, dos 6 cursos avaliados com conceito 5, 6 foram oferecidos em IES públicas, e 0, em IES privadas.
Dos 97 cursos participantes de IES privadas, o conceito 2 foi o valor modal, atribuído a 43 ( 44,3 \% ) cursos.
Em relação aos demais cursos de instituições privadas, 30 cursos ( 30,9 \% ) foram avaliados com o conceito 3, oito cursos ( 8,2 \% ) foram avaliados com o conceito 4, cinco cursos ( 5,2 \% ) foram avaliados com o conceito 1 e, nenhum curso ( 0,0 \% ), com o conceito 5.
Nessa categoria, 11 cursos ( 11,3 \% ) ficaram sem conceito ( SC ).
Na rede pública, o conceito modal foi o 3, atribuído a 19 cursos, correspondentes a 39,6 \% dos 48 cursos da categoria.
Dos demais cursos participantes, 16 cursos ( 33,3 \% ) obtiveram o conceito 4, seis cursos ( 12,5 \% ) receberam o conceito 5 e seis cursos ( 12,5 \% ), o conceito 2.
O conceito 1 foi atribuído a um curso ( 2,1 \% ), conforme mencionado.
Nessa categoria administrativa, nenhum curso ( 0,0 \% ) ficou sem conceito ( SC ).
Dos 126 cursos participantes da modalidade presencial, o conceito 3 foi o valor modal, atribuído a 46 ( 36,5 \% ) cursos.
Dos demais cursos dessa modalidade de oferta, 38 cursos ( 30,2 \% ) receberam o conceito 2 ; 24 cursos ( 19,0 \% ) receberam o conceito 4 ; seis cursos ( 4,8 \% ), o conceito 5 ; e quatro cursos ( 3,2 \% ), o conceito 1.
Nessa modalidade de oferta, oito cursos ( 6,3 \% ) ficaram sem conceito ( SC ).
Dos 19 cursos participantes da modalidade a distância, o conceito 2 foi o valor modal, atribuído a 11 ( 57,9 \% ) cursos.
Dos demais cursos \textbf{ofertados} nessa modalidade, três cursos ( 15,8 \% ) receberam conceito 3, dois cursos ( 10,5 \% ) receberam o conceito 1 ; nenhum curso ( 0,0 \% ) recebeu o conceito 5 ; e nenhum curso ( 0,0 \% ) recebeu conceito 4.
Nessa modalidade de oferta, três cursos ( 15,8 \% ) ficaram sem conceito ( SC ).
Na análise por região, observa -se, que, na região Norte, as instituições públicas participaram com cinco cursos ( 45,5 \% total regional ), dos quais nenhum ficou sem conceito ( SC ) ou recebeu o conceito 5.
O conceito modal foi 3, atribuído a três cursos, correspondendo a 60,0 \% dessa categoria na região.
O conceito 4 teve uma concentração muito próxima do modal, com um curso ( 20,0 \% ), e o conceito 2 teve um curso ( 20,0 \% ).
Os demais cursos receberam os conceitos 5 e 1, respectivamente.
As instituições privadas participaram com 6 cursos na região Norte ( 54,5 \% do total regional ), e nenhum curso ficou sem conceito ( SC ).
O conceito modal nas instituições privadas na região foi o 2, atribuído a três cursos, o que corresponde a 50,0 \% dessa categoria na região.
Os conceitos 3, 4 e 5 foram atribuídos a dois, um e nenhum curso, respectivamente.
O conceito 1 foi atribuído a nenhum curso ( 0,0 \% ).
Nenhum curso ficou sem conceito ( SC ) ou recebeu os demais conceitos.
Dos cursos oferecidos na região Norte, dez eram presenciais e nenhum ficaram sem conceito ( SC ).
Dos dez cursos oferecidos, cinco obtiveram o conceito 3 ( modal ) ; três, o conceito 2 ; dois, o conceito 4 ; nenhum, o conceito 5 ; e nenhum, o conceito 1.
Nenhum curso ficou sem conceito ( SC ) ou recebeu os demais conceitos na região Norte e na modalidade presencial.
Um curso era na modalidade a distância, tendo um recebido o conceito modal 2 ; nenhum, o conceito 5 ; nenhum, o conceito 4 ; nenhum, o conceito 3 ; e nenhum, o conceito 1.
Nenhum curso ficou sem conceito ( SC ) ou recebeu os demais conceitos na região Norte e na modalidade a distância.
Na região Nordeste, a rede privada concentra 10 dos 25 cursos participantes, o equivalente a 40,0 \% do total da região.
Entre os cursos oferecidos por IES privadas no Nordeste, cinco cursos foram avaliados com conceito 2, o conceito modal. quatro cursos obteve o conceito 3 ; um curso, o conceito 4 ; nenhum curso, o conceito 5 ; e nenhum curso obtiveram o conceito 1.
Nessa categoria administrativa e grande região, nenhum curso ficou sem conceito ( SC ).
As instituições públicas da região Nordeste participaram com 15 cursos ( 60,0 \% do total da região ).
Desses, sete ( 46,7 \% ) obtiveram conceito 3, o valor modal, seguidos por seis cursos, que receberam o conceito 4.
O conceito 5 foi atribuído a um curso ; o conceito 2, a um curso ; e o conceito 1, a nenhum curso.
Nenhum curso ficou sem conceito ( SC ).
Dos cursos oferecidos na região Nordeste, todos ( 24 ) eram presenciais : nenhum ficaram sem conceito ( SC ), nenhum obteve o conceito 1 ; cinco, o conceito 2 ; 11, o conceito 3, o valor modal ; sete, o conceito 4 ; e um, o conceito 5.
Um curso eram na modalidade a distância, tendo um curso recebido o conceito modal 2 ; nenhum, o conceito 1 ; nenhum, o conceito 3 ; nenhum, o conceito 4 ; e nenhum, o conceito 5.
Nenhum curso ficou sem conceito ( SC ) ou recebeu os demais conceitos na região Norte e na modalidade a distância.
Na região Sudeste, a proporção de cursos da rede privada, 78,8 \%, correspondeu a 52 dos 66 cursos participantes.
Nessa combinação de categoria administrativa e grande região, o conceito modal foi o 2, atribuído a 26 cursos ; e sete cursos ficaram sem conceito ( SC ).
Os demais foram avaliados com os seguintes conceitos : 3 ( 14 cursos ), 4 ( três cursos ), 1 ( dois cursos ) e 5 ( nenhum curso ).
Dos 14 cursos oferecidos por instituições públicas na região Sudeste, quatro obtiveram o conceito 5 ; dois, o conceito 4 ; cinco, o conceito 3 ; dois obtiveram conceito 2 e um obteve o conceito 1.
Nenhum curso ficou sem conceito ( SC ). atribuídos a dois, um e nenhum curso, respectivamente.
O conceito 1 foi atribuído a nenhum curso ( 0,0 \% ).
Nenhum curso ficou sem conceito ( SC ) ou recebeu os demais conceitos.
Dos cursos oferecidos na região Norte, dez eram presenciais e nenhum ficaram sem conceito ( SC ).
Dos dez cursos oferecidos, cinco obtiveram o conceito 3 ( modal ) ; três, o conceito 2 ; dois, o conceito 4 ; nenhum, o conceito 5 ; e nenhum, o conceito 1.
Nenhum curso ficou sem conceito ( SC ) ou recebeu os demais conceitos na região Norte e na modalidade presencial.
Um curso era na modalidade a distância, tendo um recebido o conceito modal 2 ; nenhum, o conceito 5 ; nenhum, o conceito 4 ; nenhum, o conceito 3 ; e nenhum, o conceito 1.
Nenhum curso ficou sem conceito ( SC ) ou recebeu os demais conceitos na região Norte e na modalidade a distância.
Na região Nordeste, a rede privada concentra 10 dos 25 cursos participantes, o equivalente a 40,0 \% do total da região.
Entre os cursos oferecidos por IES privadas no Nordeste, cinco cursos foram avaliados com conceito 2, o conceito modal. quatro cursos obteve o conceito 3 ; um curso, o conceito 4 ; nenhum curso, o conceito 5 ; e nenhum curso obtiveram o conceito 1.
Nessa categoria administrativa e grande região, nenhum curso ficou sem conceito ( SC ).
As instituições públicas da região Nordeste participaram com 15 cursos ( 60,0 \% do total da região ).
Desses, sete ( 46,7 \% ) obtiveram conceito 3, o valor modal, seguidos por seis cursos, que receberam o conceito 4.
O conceito 5 foi atribuído a um curso ; o conceito 2, a um curso ; e o conceito 1, a nenhum curso.
Nenhum curso ficou sem conceito ( SC ).
Dos cursos oferecidos na região Nordeste, todos ( 24 ) eram presenciais : nenhum ficaram sem conceito ( SC ), nenhum obteve o conceito 1 ; cinco, o conceito 2 ; 11, o conceito 3, o valor modal ; sete, o conceito 4 ; e um, o conceito 5.
Um curso eram na modalidade a distância, tendo um curso recebido o conceito modal 2 ; nenhum, o conceito 1 ; nenhum, o conceito 3 ; nenhum, o conceito 4 ; e nenhum, o conceito 5.
Nenhum curso ficou sem conceito ( SC ) ou recebeu os demais conceitos na região Norte e na modalidade a distância.
Na região Sudeste, a proporção de cursos da rede privada, 78,8 \%, correspondeu a 52 dos 66 cursos participantes.
Nessa combinação de categoria administrativa e grande região, o conceito modal foi o 2, atribuído a 26 cursos ; e sete cursos ficaram sem conceito ( SC ).
Os demais foram avaliados com os seguintes conceitos : 3 ( 14 cursos ), 4 ( três cursos ), 1 ( dois cursos ) e 5 ( nenhum curso ).
Dos 14 cursos oferecidos por instituições públicas na região Sudeste, quatro obtiveram o conceito 5 ; dois, o conceito 4 ; cinco, o conceito 3 ; dois obtiveram conceito 2 e um obteve o conceito 1.
Nenhum curso ficou sem conceito ( SC ).
Ainda na região Sudeste, grande parcela dos cursos ( 87,9 \% ) foi oferecida na modalidade presencial : 58 dos 66 cursos.
Para esses, a moda foi o conceito 2, com 23 cursos.
Os demais cursos presenciais da região foram alocados nos conceitos 3 ( 17 cursos ), 4 ( cinco cursos ), 5 ( quatro cursos ) e 1 ( três cursos ).
Seis cursos presenciais ficaram sem conceito ( SC ).
Na região Sudeste, oito cursos foram oferecidos na modalidade a distância, tendo sido nenhum com conceito 1 ; cinco com conceito 2 ; dois com conceito 3 ; nenhum com conceito 4 ; e nenhum com conceito 5.
Um curso ficou sem conceito ( SC ).
As instituições privadas concentraram 24 dos 34 cursos participantes da região Sul, 70,6 \% do total regional.
Desses, nove cursos obtiveram conceito 3, o conceito modal.
Nessa combinação de categoria administrativa e grande região, seis receberam o conceito 2 ; três receberam o conceito 4 ; dois, o conceito 1 ; e nenhum obteve o conceito 5.
Quatro cursos ficaram sem conceito ( SC ).
As instituições públicas na região Sul participaram com dez cursos ( 29,4 \% dos cursos da região ), aos quais foram atribuídos os conceitos 1 ( nenhum curso ), 2 ( nenhum curso ), 3 ( três cursos ), 4 ( seis cursos ) e 5 ( um curso ).
Nenhum curso da região Sul oferecido por IES públicas ficou sem conceito ( SC ).
Na região Sul, 26 dos cursos foram oferecidos na modalidade presencial.
Para esses, a moda foi o conceito 3, com 11 cursos.
Os demais cursos presenciais da região foram alocados nos conceitos 4 ( nove cursos ), 2 ( três cursos ), 5 ( um curso ) e 1 ( nenhum curso ).
Na modalidade a distância, foram oferecidos oito cursos, sendo dois cursos com conceito 1, três cursos com conceito 2, um curso com conceito 3 ( modal ), nenhum curso com conceito 4 e nenhum curso com conceito 5.
Dois cursos ficaram sem conceito ( SC ) nessa região.
Na região Centro-Oeste, 5 dos 9 cursos participantes eram de IES privadas ( 55,6 \% em termos regionais ).
Desses, três cursos receberam conceito 2, o valor modal.
Dos demais cursos, um recebeu o conceito 3, um receberam o conceito 1, nenhum receberam o conceito 5 e nenhum recebeu o conceito 4. nenhum ficou sem conceito ( SC ).
Quanto aos quatro cursos oferecidos por instituições públicas na região Centro-Oeste ( 44,4 \% do total regional ), nenhum ficou com conceito 1 ; dois, com conceito 2 ; um, com conceito 3 ; um, com conceito 4 ; e nenhum ficou com conceito 5.
Nenhum ficou sem conceito ( SC ) nessa região.
Também na região Centro-Oeste, a maioria dos cursos foi \textbf{ofertada} na modalidade presencial : 8 ( 88,9 \% ).
Para esses, a moda foi o conceito 2, com 4 cursos ( 50,0 \% ).
Os demais cursos presenciais da região foram alocados nos conceitos 3 ( dois cursos, 25,0 \% ), 4 ( um curso, 12,5 \% ), 1 ( um curso, 12,5 \% ) e 5 ( nenhum curso, 0,0 \% ).
Nenhum curso presencial da região ficou sem conceito ( SC ).
Na modalidade a distância, foi oferecido um curso, sendo : nenhum curso com conceito 1, um curso com conceito 2, nenhum curso com conceito 3 ( modal ), nenhum curso com conceito 4 e nenhum curso com conceito 5.
Nenhum curso ficou sem conceito ( SC ) nessa região.
Para o único curso a distância nessa região, o conceito atribuído foi 3.
A título de síntese, no Gráfico 4.2, apresenta -se a distribuição do Conceito Enade, segundo a categoria administrativa da IES.
Nesse gráfico, \textbf{poligonais} mais à esquerda representam Conceito Enade mais baixos, e \textbf{poligonais} mais à direita, Conceito Enade mais altos.
É importante observar que, para cada \textbf{poligonal} presente nesse gráfico, a porcentagem constante no eixo y refere -se ao universo de cursos oferecidos por instituições públicas ( \textbf{poligonal} azul ) ou privadas ( \textbf{poligonal} verde ), respectivamente.
Os cursos em IES públicas ( \textbf{poligonal} azul ) apresentam uma \textbf{poligonal} mais à direita do que os em IES privadas ( \textbf{poligonal} verde ) e, consequentemente, uma distribuição de Conceito Enade com valores maiores. 4.3 CONCEITOS POR ORGANIZAÇÃO ACADÊMICA E GRANDE REGIÃO Na Tabela 4.3, encontra -se a distribuição dos conceitos atribuídos aos cursos participantes do Enade 2023, na área de Engenharia de Computação I, por organização acadêmica, segundo as grandes regiões brasileiras.
Como mostram os dados da Tabela 4.3, dos 145 cursos de Engenharia de Computação I participantes, 78 eram oferecidos em universidades, 38 em centros universitários e 21 em faculdades ; oito cursos foi oferecido por CEFET/IF.
Essa distribuição corresponde a, respectivamente, 53,8 \%, 26,2 \%, 14,5 \% e 5,5 \% dos cursos.
De acordo com os dados apresentados na mesma tabela, de todos os 49 cursos avaliados com o conceito 3, 25 eram vinculados a universidades.
Os cursos vinculados a esse tipo de organização acadêmica tiveram o conceito 2 como conceito modal, atribuído a 26 cursos ( 33,3 \% ).
Os demais cursos \textbf{ofertados} por universidades receberam os conceitos 4 ( 19 cursos ), 5 ( cinco cursos ) e 1 ( um curso ).
Dois cursos ficaram sem conceito ( SC ).
Dos cursos vinculados a centros universitários, o conceito modal também foi o 2, atribuído a 16 cursos ( 42,1 \% ).
Seis cursos vinculados a esse tipo de organização acadêmica ficaram sem conceito ( SC ), e os outros cursos receberam os conceitos 3 ( dez cursos ), 4 ( três cursos ), 1 ( três cursos ) e 5 ( nenhum curso ).
Dos 21 cursos mantidos por faculdades, 7 ( 33,3 \% ) receberam conceito 3, o conceito modal.
O segundo conceito mais frequentemente atribuído aos cursos oferecidos por essa categoria de IES foi o conceito 2, recebido por seis cursos.
Na sequência, dois cursos obtiveram conceito 4 ; dois cursos, o conceito 1 ; e um curso, o conceito 5.
Dos cursos oferecidos em faculdades, três cursos ficaram sem conceito ( SC ).
Dos cursos oferecidos por CEFET/IF, o conceito modal foi o 3, atribuído a sete cursos ( 87,5 \% ).
Nenhum curso vinculado a esse tipo de organização acadêmica ficou sem conceito ( SC ), e os outros cursos receberam os conceitos 2 ( um curso ), 5 ( nenhum curso ), 4 ( nenhum curso ) e 1 ( nenhum curso ), respectivamente.
Considerando -se, separadamente, as regiões brasileiras, verifica -se que, na região Norte, as universidades participaram com seis cursos avaliados, aos quais foram atribuídos os seguintes conceitos : 1 a nenhum curso ; 2 a dois cursos ; 3, o conceito modal, a três cursos ; conceito 4 a um curso e 5 a nenhum curso.
Nenhum curso ficou sem conceito ( SC ).
Os centros universitários da região Norte foram representados por dois cursos, dos quais nenhum curso obteve conceito 1 ; um curso, conceito 2 ; nenhum curso, conceito 3, o valor modal ; um curso, o conceito 4 e nenhum curso, o conceito 5.
Nenhum curso ficou sem conceito ( SC ).
As faculdades participaram com três cursos na região Norte, os quais receberam os conceitos 1 ( nenhum curso ), 2 ( um curso ), 3 ( dois cursos ), 4 ( nenhum curso ) e 5 ( nenhum curso ).
Nenhum curso ficou sem conceito ( SC ).
Os CEFET/IF participaram com nenhum curso, aos quais foram atribuídos os seguintes conceitos : 1 a nenhum curso ; 2 a nenhum curso ; 3 a nenhum curso ; 4 a nenhum curso ; e 5 a nenhum curso.
Nenhum curso ficou sem conceito ( SC ).
Na região Nordeste, as universidades participaram com 18 dos 25 cursos da área de Engenharia de Computação I oferecidos.
O conceito modal foi 3, atribuído a sete cursos.
Outros seis cursos obtiveram o conceito 4 ; quatro cursos, o conceito 2 ; um curso, o conceito 5 ; e nenhum curso, o conceito 1.
Nenhum curso ficou sem conceito ( SC ).
Os centros universitários \textbf{contaram} com três cursos participantes na região Nordeste, dois dos quais receberam o conceito modal 2.
Para os demais cursos, foram atribuídos os conceitos 4 ( um curso ), 5 ( nenhum curso ), 3 ( nenhum curso ) e 1 ( nenhum curso ).
Nenhum curso oferecido por centros universitários no Nordeste ficaram sem conceito ( SC ).
As faculdades foram representadas por um curso na região Nordeste, os quais receberam os conceitos 1 ( nenhum curso ), 2 ( nenhum curso ), 3 ( um curso ), 4 ( nenhum curso ) e 5 ( nenhum curso ).
Nenhum curso ficou sem conceito ( SC ).
Na região Nordeste, os CEFET/IF participaram com três cursos, aos quais foram atribuídos os seguintes conceitos : 1 a nenhum curso ; 2 a nenhum curso ; 3 a três cursos ; 4 a nenhum curso ; e 5 a nenhum curso.
Nenhum curso ficou sem conceito ( SC ).
Na região Sudeste, as universidades concentraram 30 dos 66 cursos de Engenharia de Computação I da região.
Dos cursos oferecidos em IES com este tipo de organização acadêmica, nessa região, o conceito modal foi 2, atribuído a 14 cursos.
Os demais cursos receberam os conceitos 3 ( oito cursos ), 5 ( três cursos ), 4 ( dois cursos ) e 1 ( um curso ).
Dois cursos ficaram sem conceito ( SC ).
Os centros universitários participaram com 21 cursos na região Sudeste, dos quais 11 obtiveram o conceito modal 2.
O conceito 3 foi atribuído a sete cursos ; o conceito 4 a um curso ; o conceito 5 a nenhum curso ; e o conceito 1 a nenhum curso.
Dois cursos ficaram sem conceito ( SC ) nessa combinação de organização acadêmica e região.
Já as faculdades foram representadas por 13 cursos na região Sudeste, dos quais três ficaram sem conceito ( SC ).
O conceito modal 2 foi atribuído a 3 cursos.
Aos demais cursos foram atribuídos os conceitos 4 ( dois cursos ), 3 ( dois cursos ), 1 ( dois cursos ) e 5 ( um curso ).
Os CEFET/IF participaram com dois cursos na região Sudeste, aos quais foram atribuídos os seguintes conceitos : 1 a nenhum curso ; 2 a nenhum curso ; 3 a dois cursos ; 4 a nenhum curso ; e 5 a nenhum curso.
Nenhum curso ficou sem conceito ( SC ).
Dos 34 cursos da região Sul, 20 eram mantidos por universidades, para os quais o conceito modal foi 4, atribuído a 9 cursos.
Os demais cursos receberam os conceitos 3 ( sete cursos ), 2 ( três cursos ), 5 ( um curso ) e 1 ( nenhum curso ).
Em IES com esse tipo de organização acadêmica, nenhum curso da região Sul ficou sem conceito ( SC ).
Dos dez cursos participantes de centros universitários da região Sul, quatro ficaram sem conceito ( SC ), e dois cursos receberam conceito 3, o valor modal.
Os conceitos \textbf{restantes} ficaram assim distribuídos : 2 ( dois cursos ), 1 ( dois cursos ), 5 ( nenhum curso ) e 4 ( nenhum curso ).
Quanto aos três cursos vinculados a faculdades na região Sul, dois cursos receberam o conceito 3, valor modal.
Os demais conceitos ficaram assim distribuídos : 2 ( um curso ), 5 ( nenhum curso ), 4 ( nenhum curso ) e 1 ( nenhum curso ).
Dos três cursos, nenhum ficou sem conceito ( SC ).
Os CEFET/IF participaram com um curso na região Sul, aos quais foram atribuídos os seguintes conceitos : 1 a nenhum curso ; 2 a nenhum curso ; 3 a um curso ; 4 a nenhum curso ; e 5 a nenhum curso.
Nenhum curso ficou sem conceito ( SC ).
Na região Centro-Oeste, 4 dos 9 cursos oferecidos eram mantidos por universidades.
Para esse tipo de organização acadêmica, o conceito modal foi 2, atribuído a três cursos.
Em relação às universidades, nenhum curso da região ficou sem conceito ( SC ).
Os demais cursos receberam os conceitos 4 ( um curso ), 5 ( nenhum curso ), 3 ( nenhum curso ) e 1 ( nenhum curso ).
Os centros universitários da região Centro-Oeste \textbf{contaram} com dois cursos : um curso foi avaliado com o conceito 3 ( modal ) ; um curso, com o conceito 1 ; nenhum curso, com o conceito 5 ; nenhum curso, com o conceito 4 ; e nenhum curso obteve o conceito 2.
Nenhum curso ficou sem conceito ( SC ).
O único curso oferecido em faculdade na região Centro-Oeste, recebeu o conceito 2.
Nessa região, nenhum curso oferecido em faculdades ficou sem conceito ( SC ) ou obteve os conceitos 1 e 5.
No que se refere aos CEFET/IF da região Centro-Oeste, dois cursos foram avaliados : um curso deles com o conceito 3 ( modal ) ; um curso com o conceito 2 ; nenhum curso com o conceito 5 ; nenhum curso com o conceito 4 ; nenhum curso obteve o conceito 1 ; e nenhum curso ficou sem conceito ( SC ).
De maneira semelhante à forma como se \textbf{procedeu} em relação à seção anterior, esta seção encerra -se com o Gráfico 4.3, que apresenta a distribuição do Conceito Enade dos cursos de Engenharia de Computação I segundo a organização acadêmica da IES.
No Gráfico 4.3, apresenta -se a distribuição do Conceito Enade segundo a organização acadêmica da IES.
A leitura desse último gráfico é semelhante à dos anteriores, seguindo -se a lógica de que \textbf{poligonais} deslocadas para a esquerda correspondem aos conceitos mais baixos ( 1, 2 e 3 ) e \textbf{poligonais} deslocadas para a direita correspondem a Conceitos Enade mais altos.
Nas universidades ( linha azul ), todos os cursos obtiveram conceito 5 ou abaixo disso, já que a \textbf{poligonal} \textbf{cumulativa} alcança 100 \% para o conceito 5.
Além disso, os cursos de universidades apresentam uma \textbf{poligonal} mais à direita do que as demais, \textbf{denotando} uma melhor distribuição de Conceitos Enade.
No outro extremo, encontram -se os cursos mantidos por centros universitários ( linha verde ), apresentando as respectivas \textbf{poligonais} mais à esquerda e, consequentemente, uma distribuição de Conceitos Enade com os menores valores.
Os cursos mantidos por universidades ( linha vermelha ) apresentam \textbf{poligonais} mais centrais e, consequentemente, uma distribuição de Conceitos Enade com os valores intermediários.

% matched lemmas: abscissa, contar, cruzar, cumulativo, denotar, equivaler, finalidade, indicado, ministrar, ofertar, ordenado, panorama, poligonal, pontilhado, proceder, reiterativo, restante, roxa, segunda, tabular, tipologia, variável, visualizar
\end{textsample}
